\documentclass{article}
\usepackage[utf8]{inputenc}
\usepackage[spanish]{babel}
\usepackage{geometry}
\usepackage{array}
\usepackage{amsmath}

\geometry{letterpaper, margin=2.5cm}

\begin{document}

Nombre: Rivas Bodero Jordy \\
EV-PRAC-01-SUM-Procesos y el Framework de IR. Resolver el problema

\section*{D1. Comprensión Conceptual}
\textbf{T1.} Requisitos, docuementos StRS/SyRS/SRS, y actividades del progreso. \\
- Un requisito es una necesidad, capacidad o condición que debe cumplir un sistema, producto o servicio para satisfacer una necesidad de un necesitado /interesado, cumplir un objetivo o ajustarse a una condición establecida.

\subsection*{b. Diferencias entre StRS, SyRS y SRS}
\begin{center}
\begin{tabular}{|p{2cm}|p{4cm}|p{1.5cm}|p{5cm}|}
\hline
\textbf{Documento} & \textbf{¿Qué describe?} & \textbf{Ejemplo} & \textbf{Justificación} \\ \hline
StRS & Necesidades y expectativas de los interesados. & $E_1$ & La coordinadora necesita recibir alertas inmediatas. \\ \hline
SyRS & Requisitos del sistema completo incluyendo sus componentes e interacciones. & $E_6$ & Establece que el sistema debe continuar registrando temperaturas durante cortes electricos\\ \hline
SRS & Requisitos específicos del software y sus funcionalidades. & $E_5$ & Establece que el software genera informes mensuales en PDF. \\ \hline
\end{tabular}
\end{center}

Diferencias que debe recordar:
\begin{itemize}
    \item StRS = Qué necesitan los interesados.
    \item SyRS = Qué debe cumplir el sistema completo.
    \item SRS = Qué debe cumplir el software.
\end{itemize}

\subsection*{c. Actividades del proceso}
\begin{center}
\begin{tabular}{|p{1cm}|p{3.5cm}|p{2cm}|p{6cm}|}
\hline
\textbf{ID} & \textbf{Actividad de IR} & \textbf{Tipo} & \textbf{Justificación} \\ \hline
$H_1$ & Elicitación & Central & Se realizan entrevistas para obtener las necesidades de los interesados. \\ \hline
$H_2$ & Análisis y negociación & Central & Se comparan las propuestas y correos y se llega a un acuerdo. \\ \hline
$H_3$ & Especificación & Central & Se documentan los requisitos acordados con identificadores unicos. \\ \hline
$H_4$ & Validación & Central & Los interesados revisan el documento y detectan requisitos incompletos. \\ \hline
$H_5$ & Gestión & Transversal & Se controla un cambios, se analizan impactos y se publica una nueva versión. \\ \hline
\end{tabular}
\end{center}

\section*{D2. Aplicación y análisis}
\textbf{T2. (a)} Clasificación de requisito.
\begin{center}
\begin{tabular}{|p{1cm}|p{2cm}|p{2cm}|p{5.5cm}|}
\hline
\textbf{ID} & \textbf{Tipo} & \textbf{Documento} & \textbf{Justificación} \\ \hline
$E_1$ & Funcional & StRS & Expresa la necesidad de recibir alertas ante excursiones términicas \\ \hline
$E_2$ & Dominio & SyRS & Establece el rango que se deben conservar las vacunas \\ \hline
$E_3$ & Restricción & SyRS & Establece el protocolo a seguir y la frecuencia \\ \hline
$E_4$ & Restricción del proyecto & Plan de proyecto & Fija una fecha límite \\ \hline
$E_5$ & Funcional & SRS & Genera informes en PDF \\ \hline
$E_6$ & No Funcional & SyRS & La continuidad durante los cortes de luz \\ \hline
$E_7$ & No es un requisito del producto & Contexto del sistema actual & Registro manual de temperatura(descripción) \\ \hline
$E_8$ & No Funcional & SRS & Tiempo de respuesta con 50 usuarios.\\ \hline
$E_9$ & Restricción & SyRS & Limita el entorno del Sistema operativo \\ \hline
$E_{10}$ & Funcional & SRS & Registro de permisos a gente autorizada \\ \hline
$E_{11}$ & No Funcional & SRS & Impone una condición medible \\ \hline
$E_{12}$ & Requisito del proyecto & Plan de proyecto & Define como se establecen los objetivos. \\ \hline
\end{tabular}
\end{center}

\subsection*{(b) Autoría de 6 requisitos}
\begin{center}
\begin{tabular}{|p{1cm}|p{6cm}|p{4cm}|}
\hline
\textbf{ID} & \textbf{Termino o construcción de la problemática}&\textbf{Características} \\ \hline
$B_1$ & Fácil de usar y rápido  & No ambiguo y verificable \\ \hline
$B_2$ & Alertas oportunas y temperatura inadecuada  & No ambiguo y verificable  \\ \hline
$B_3$ & Sin conexión.  & Factible y verificable \\ \hline
$B_4$ & Alojamiento y gestor de datos  & Apropiado \\ \hline
$B_5$ & Exportar (Excel, cs, etc)  & No ambiguo y verificable \\ \hline
$B_6$ & Ninguno  & Cumple las caracteristicas \\ \hline
\end{tabular}
\end{center}

\section*{D3. Síntesis y producción}
\textbf{T3.} Redacción de seis requisitos Bien Formados.
\begin{center}
\begin{tabular}{|p{1cm}|p{10cm}|}
\hline
\textbf{ID} & \textbf{Requisito} \\ \hline
$B_1$ & Cuando existan 50 usuarios conectados simultáneamente, el sistema debera mostrar el historial de temperaturade un establecimiento en menos de 3 segundos. \\ \hline
$B_2$ & Cuando la temperatura esté fuera de los rangos establecidos se enviaran notificaciones al responsable del establecimiento en un plazo de 5 minutos. \\ \hline
$B_3$ & El sistema debe recibir y almacenar las lecturas de temperatura registradas por los sensores dentro del tiempo establecidos \\ \hline
$B_5$ & Cuando el sistema reciba una solicitud de importar datos este tendra establecido el formato al cual se solicite se proporcionara el respectivo archivo. \\ \hline
$B_6$ & El sistema enviara notificaciones al responsable en caso de que existan anomalías en la temperatura. \\ \hline
\end{tabular}
\end{center}

\section*{D4. Uso de estándares y referencias}
\textbf{T4.} y las justificaciones de $T_1$ a $T_3$
\begin{center}
\begin{tabular}{|p{2cm}|p{5cm}|p{3cm}|}
\hline
Específico & Define que debe hacer el sistema, sin ambiguedades & No ambiguo \\ \hline
Medible & Permite medir el objetivo & Verificable \\ \hline
Alcanzable & Al alcance de los involucrados & Factible \\ \hline
Realizable & Dentro de las capacidades & Factible \\ \hline
Trazable & Va en contexto con lo solicitado & Identificable \\ \hline
\end{tabular}
\end{center}

\subsection*{Defectos}
\begin{itemize}
    \item No existe completitud en el conjunto
    \item Inconsistencia entre los requisitos
    \item Conflicto en la continuidad del requisito
    \item Diferencia entre la lectura inmediata y periódica
    \item Exportación de datos sin definir alcance
    \item Falta de uniformidad entre las alertas
\end{itemize}

\begin{center}
\begin{tabular}{|p{3cm}|p{1cm}|p{4.5cm}|p{4.5cm}|}
\hline
\textbf{Categorización} & \textbf{ID} & \textbf{Defecto} & \textbf{Corrección} \\ \hline
Consistencia & $E_6, B_3$ & $E_6$ establece continuidad durante cortes de hasta 8 horas; $B_3$ exige registro indefinido incluso sin conexión. & Definir límite de autonomía y comportamiento ante la pérdida de energía. \\ \hline
Completitud & $E_2, B_6$ & $E_2$ define el rango de $+2^{\circ}\text{C}$ a $+8^{\circ}\text{C}$, pero $B_6$ establece alertas cuando se supera $+8^{\circ}\text{C}$. & Añadir la condición de alerta cuando la temperatura sea inferior a $+2^{\circ}\text{C}$. \\ \hline
\end{tabular}
\end{center}

\end{document}